% EdgeSched: phase-aware elastic residency and joint resource scheduling.
% This standalone note summarizes the mathematical model used by the design.

\documentclass[10pt]{article}
\usepackage[T1]{fontenc}
\usepackage{newtxtext,newtxmath}
\usepackage[margin=0.72in]{geometry}
\usepackage{amsmath,amsfonts}
\usepackage{booktabs,array}
\usepackage{xcolor}

\newcommand{\CPU}{\textsc{cpu}}
\newcommand{\GPU}{\textsc{gpu}}
\newcommand{\Unmanaged}{\textsc{unmanaged}}

\begin{document}

\begin{center}
{\LARGE\bfseries EdgeSched: Phase-Aware Elastic Residency\par}
\vspace{2pt}
{\Large\bfseries and Joint Resource Scheduling\par}
\vspace{5pt}
{\small Design formulation and current workload instantiation}
\end{center}

\section*{1. Phase Capability and Measured Profiles}

EdgeSched schedules an executable \emph{phase}.  A phase may be an LLM prefill
or decode, a TTS generator or codec, a diffusion job, or an auxiliary stage.
Its static capability declaration is

\begin{equation}
\Phi_i = \left\langle
    d_i,\rho_i,g_i,\chi_i,
    \operatorname{role}_i,
    \operatorname{preempt}_i,
    \operatorname{queue}_i,
    \operatorname{deps}_i
\right\rangle .
\end{equation}

Here, $d_i$ is the device affinity; $\rho_i$ is fixed, whole-model migratable,
or partitionable residency; and

\begin{equation}
g_i\in\{\mathrm{NONE},\mathrm{WHOLE},\mathrm{LAYER},
             \mathrm{COMPONENT},\mathrm{STATE}\}
\end{equation}

records the validated reconfiguration granularity.  The value $\mathrm{NONE}$
applies to a fixed phase and $\mathrm{WHOLE}$ to a phase that can only move as a
complete instance.  GPU compute-share capability is

\begin{equation}
\chi_i\in\{\mathrm{UNMANAGED},\mathrm{CAP\_LIMITABLE},
                 \mathrm{SM\_PARTITIONABLE}\}.
\end{equation}

The remaining fields define service semantics, preemption and queueing rules,
and workflow dependencies.  A declared capability does not by itself create a
scheduling option: only configurations supported by the backend and validated
by profiling enter the finite candidate set $\mathcal U_i$.

Let $u_i=(r_i,c_i,s_i,b_i)\in\mathcal U_i$ denote a residency profile, physical
CPU-core budget, GPU compute-share setting, and batch/concurrency setting.  For
an unmanaged GPU client, $s_i=\bot$ rather than a fictitious full-device
reservation.  Offline profiling records

\begin{equation}
\Psi_i(u_i\mid z_i,h_t)=\left\langle
    \widehat M_i^{G},\widehat M_i^{H},
    \widehat T_i^{\mathrm{first}},\widehat T_i^{\mathrm{steady}},
    \widehat Q_i
\right\rangle ,
\end{equation}

where $z_i$ contains input size, context/KV size, model and runtime settings,
and $h_t$ contains recent observations used to update the offline profile.
$\widehat Q_i$ contains workload-specific predictions such as TTFT,
sentence-ready latency, RTF, token gap, throughput, and slowdown.  Unsupported
or unverified configurations are absent rather than assigned optimistic values.

Transitions are profiled separately because they are edges between two
residency configurations, not properties of one configuration:

\begin{align}
\Gamma_i(r_i\!\rightarrow r_i'\mid K_i,z_i)
   &=\left\langle T_i^{\mathrm{load}},T_i^{\mathrm{save}},
      T_i^{\mathrm{restore}},T_i^{\mathrm{cutover}},
      T_i^{\mathrm{cleanup}}\right\rangle ,\\
C_i^{\mathrm{tr}}(r_i\!\rightarrow r_i'\mid K_i,z_i)
   &=T_i^{\mathrm{load}}+T_i^{\mathrm{save}}
     +T_i^{\mathrm{restore}}+T_i^{\mathrm{cutover}}
     +T_i^{\mathrm{cleanup}}.
\end{align}

\section*{2. Current Workload Instantiation}

\begin{center}
\footnotesize
\setlength{\tabcolsep}{3pt}
\begin{tabular}{>{\raggedright\arraybackslash}p{0.10\linewidth}
                >{\raggedright\arraybackslash}p{0.17\linewidth}
                >{\raggedright\arraybackslash}p{0.25\linewidth}
                >{\raggedright\arraybackslash}p{0.19\linewidth}
                >{\raggedright\arraybackslash}p{0.20\linewidth}}
\toprule
Phase & Service semantics & Validated residency candidates
      & GPU-control evidence & Scheduling obligation \\
\midrule
LLM & resident, stateful, interactive
    & CPU, profiled layer-hybrid configurations, and full GPU
    & unmanaged measured; controlled profiles pending platform validation
    & remain reachable and preserve context across transitions \\
TTS & burst or session-warm, deadline-sensitive
    & CPU fallback, T0 full GPU, and T3 component-hybrid
    & unmanaged measured; controlled profiles pending platform validation
    & minimize sentence-ready delay and meet playback deadlines \\
Diffusion & burst, queueable, GPU-required
    & full GPU execution
    & unmanaged measured; controlled profiles pending platform validation
    & admit only when capacity-safe; do not interrupt an admitted job \\
Auxiliary & short helper phases
    & measured CPU-preferred configurations
    & not required in the current plan
    & avoid starving system, audio, and interactive threads \\
\bottomrule
\end{tabular}
\end{center}

The controlled GPU-sharing code path is implemented, but a controlled profile
does not enter $\mathcal U_i$ until the target platform passes capability and
performance validation.  The table is a workload instantiation of $\Phi_i$;
the scheduler interface is not limited to these model names.

\section*{3. Runtime State and Event-Driven Actions}

At a scheduling event, EdgeSched observes

\begin{equation}
x_t=\left\langle
\mathbf q_t,\mathbf K_t,
{\bigl\{r_i^{\mathrm{cur}}\bigr\}}_{i\in\mathcal A_t},
M_G,M_H,S_G,U_C,\mathbf d_t,v_t,h_t
\right\rangle .
\end{equation}

$\mathbf q_t$ contains per-phase queues and running states; $\mathbf K_t$
contains phase state sizes such as LLM context and KV; $\mathbf d_t$ contains
deadlines; and $h_t$ contains recent durations, arrivals, interference, and
transitions.  In the current workload, $\mathbf q_t$ includes the diffusion-job
queue and TTS-chunk queue.  Decisions occur at arrivals, completions, queue or
session changes, completed reconfigurations, resource-boundary violations, and
backend failures---not after every token or utilization fluctuation.

The current workload uses three event modes:

\begin{align*}
S_0 &: \text{no GPU-required burst is admitted: optimize the interactive pipeline},\\
S_1 &: \text{a GPU-required burst is running: protect the job and interaction SLOs},\\
S_2 &: \text{the burst queue is nonempty: avoid restore--evict oscillation}.
\end{align*}

Diffusion instantiates the GPU-required burst in the present experiments; the
mode definitions also apply to another profiled burst phase with the same
service obligations.

The joint action is

\begin{equation}
a_t=\left({\bigl\{u_i=(r_i,c_i,s_i,b_i)\bigr\}}_{i\in\mathcal A_t},\mu_t\right),
\end{equation}

where $\mu_t$ is hold, urgent release, atomic restore, or reconfiguration.

\section*{4. Hierarchical Feasibility}

EdgeSched first filters residency and capacity, and only then considers compute
sharing.  A candidate belongs to $\mathcal F_{\mathrm{cap}}(x_t)$ only if

\begin{align}
\sum_{i\in\mathcal A_t}
  \widehat M_i^G(u_i\mid z_i,h_t)+G_{\mathrm{safe}}
  &\le V_{\mathrm{GPU}},\\
\sum_{i\in\mathcal A_t}
  \widehat M_i^H(u_i\mid z_i,h_t)+H_{\mathrm{safe}}
  &\le V_{\mathrm{host}},\\
\sum_{i\in\mathcal C_t}c_i+R_{\mathrm{sys}}
  &\le C_{\mathrm{physical}}.
\end{align}

$\mathcal C_t$ contains phases receiving physical CPU cores, while
$R_{\mathrm{sys}}$ protects system, scheduler, sampling, and audio work.  For a
backend that provides verified disjoint SM partitions,

\begin{equation}
\sum_{i\in\mathcal G_t}s_i\le S_{\mathrm{GPU}},
\qquad s_i\in\mathcal S_i .
\end{equation}

For an MPS active-thread cap, $s_i$ is a verified client limit rather than an
exclusive reservation, so the disjoint-partition interpretation is not used.
For unmanaged execution $s_i=\bot$, and feasibility relies on measured
co-running profiles rather than an invented SM allocation.

Capacity feasibility is necessary but not sufficient.  The SLO-feasible set is

\begin{equation}
\mathcal F_{\mathrm{SLO}}(x_t)=
\left\{a\in\mathcal F_{\mathrm{cap}}(x_t):
g_k(a\mid x_t)\le 0,\ \forall k\in\mathcal D_t\right\},
\end{equation}

where $g_k$ represents a frozen interaction, playback, slowdown, correctness,
or backend-validity constraint.  Thus VRAM-compatible co-residency does not
imply execution-compatible co-scheduling.

\section*{5. Cost-Aware Pareto Planning}

For already-arrived work over decision horizon $H$, configuration-specific
predictors estimate phase duration, queue completion, interference, and the
critical path of the workflow.  EdgeSched does not wait for hypothetical future
requests in order to create a batch.  Within $\mathcal F_{\mathrm{SLO}}(x_t)$,
the running-performance vector is

\begin{equation}
\mathbf J_{\mathrm{run}}(a\mid x_t)=\left\langle
\widehat T_{\mathrm{interactive}},
\widehat T_{\mathrm{burst}},
\widehat T_{\mathrm{workflow}},
-\widehat\Theta_{\mathrm{interactive}},
N_{\mathrm{reconfig}}
\right\rangle .
\end{equation}

Only components relevant to the active workflow are instantiated.  The planner
first constructs the non-dominated set

\begin{equation}
\mathcal P(x_t)=
\operatorname*{ParetoMin}_{a\in\mathcal F_{\mathrm{SLO}}(x_t)}
\mathbf J_{\mathrm{run}}(a\mid x_t),
\end{equation}

then applies a fixed, pre-declared lexicographic policy to select

\begin{equation}
a_t^{\mathrm{cand}}=
\operatorname*{LexMin}_{a\in\mathcal P(x_t)}
\mathbf J_{\mathrm{run}}(a\mid x_t).
\end{equation}

This avoids an opaque weighted sum of latency, throughput, and SLO penalties
with incompatible units.  The priority order and all SLO thresholds are frozen
before formal evaluation.

For the current non-incremental TTS backend, the measured first-output metric is
the time at which the complete first-sentence waveform becomes available:

\begin{equation}
T_{\mathrm{sentence\ ready}}=t_1^{\mathrm{ready}}-t_{\mathrm{request}}.
\end{equation}

When transition, LLM generation, and TTS execution are serialized, this equals
their sum; when they overlap, EdgeSched measures the actual critical path rather
than adding overlapping intervals.  Playback feasibility requires

\begin{equation}
t_{j+1}^{\mathrm{ready}}\le t_j^{\mathrm{playback\ end}},
\qquad j\ge 1.
\end{equation}

The metric may be renamed first-audio latency only after an incremental audio
backend exposes the first playable samples.

\section*{6. Hysteresis and Safe Reconfiguration}

Let $a_{\mathrm{cur}}$ be the current action and $a'$ the selected feasible
candidate.  Define the predicted cumulative running-latency improvement over
already-arrived work, excluding the one-time transition itself, as

\begin{equation}
\Delta_H(a_{\mathrm{cur}},a'\mid x_t)=
\widehat L_H(a_{\mathrm{cur}}\mid x_t)
-\widehat L_H(a'\mid x_t).
\end{equation}

EdgeSched switches only when

\begin{equation}
\Delta_H(a_{\mathrm{cur}},a'\mid x_t)
> C^{\mathrm{tr}}(a_{\mathrm{cur}}\!\rightarrow a'\mid x_t)+\eta,
\qquad t-t_{\mathrm{last\ switch}}\ge\tau_{\min},
\end{equation}

where $C^{\mathrm{tr}}$ sums the transition-edge costs of changed phases,
$\eta$ is a hysteresis margin, and $\tau_{\min}$ is a minimum residency window.
During a burst train, the known queue extends horizon $H$, allowing EdgeSched to
retain a degraded or hybrid residency configuration until the queue drains
instead of restoring after every job.

\end{document}
